\section{Entropy of twisted closures}%
\label{sec:peps_twisted_closure_entropy}

Group elements inserted on the regular bonds of a region twist its open
tensor. This section computes the cross Gram matrices of twisted regions,
shows that closures with complementary twists have a common reduced density,
and deduces the entropy of superpositions of closures on rectangles inside the
torus.

\begin{lemma}[Expansion of a twisted-region cross Gram matrix]%
    \label{thm:peps_regular_twisted_gram_expansion}
    \lean{TNLean.PEPS.regularOpenRegionMatrix_apply_group,
        TNLean.PEPS.regularOpenRegionMatrix_crossGram_apply_group,
        TNLean.PEPS.prod_twistedSiteGram_eq_sum_compatible,
        TNLean.PEPS.regularTwistedOpenRegionMatrix_crossGram_eq_sum_compatible}
    \leanok
    \uses{def:peps_regular_twisted_region}
    Suppose the original site maps satisfy $A_v^\dagger A_v=c_v\Pi_v$.
    The entries of $(T_R^u)^\dagger T_R^w$ are obtained by summing over
    two internal-bond configurations and one translation $q_v$ per vertex.
    At each incident half-bond the constraint is
    \begin{align}
        \eta^u_{v,e}=q_v\theta^w_{v,e}.
    \end{align}
    The prescribed crossing configurations are retained, and the scalar
    from the local averages is $(\prod_{v\in R}c_v)|G|^{-|R|}$.
\end{lemma}

\begin{proof}\leanok
    \uses{thm:peps_regular_site_gram}
    Expanding both open contractions, summing the physical indices
    independently, and substituting each original local Gram identity at its
    two transformed virtual configurations gives the result.
\end{proof}

\begin{lemma}[Synchronization through common identity bonds]%
    \label{thm:peps_regular_twisted_synchronization}
    \lean{TNLean.PEPS.regionLabels_eq_of_subgraph_internalEdge_eq,
        TNLean.PEPS.IsTwistedRegionLabelCompatible,
        TNLean.PEPS.IsTwistedRegionLabelCompatible.tail_eq,
        TNLean.PEPS.IsTwistedRegionLabelCompatible.head_eq,
        TNLean.PEPS.IsTwistedRegionLabelCompatible.internalEdge_intertwining,
        TNLean.PEPS.IsTwistedRegionLabelCompatible.exists_common_label,
        TNLean.PEPS.isTwistedRegionLabelCompatible_iff_exists_translation,
        TNLean.PEPS.sum_twistedRegionLabelCompatible_eq_sum_translation}
    \leanok
    Suppose the region has a connected spanning subgraph on whose bonds
    both assignments $u,w$ are the identity. Suppose also that both are
    the identity on crossing bonds. Then compatibility is equivalent to
    one common translation $x$, with
    \begin{align}
        q_v=x,\qquad \eta_e=x\theta_e,\qquad
        u_ex=xw_e\quad\text{for every internal bond }e.
    \end{align}
    Consequently the sum over compatible vertex translations is the
    corresponding restricted sum over $x\in G$.
\end{lemma}

\begin{proof}\leanok
    \uses{thm:peps_regular_twisted_gram_expansion}
    The tail constraint gives $\eta_e=q_{s(e)}\theta_e$, while cancellation
    in the head constraint gives $u_eq_{s(e)}=q_{t(e)}w_e$.
    Along each common identity bond this equates the endpoint translations,
    so connectedness makes them constant, and the boundary identity operators then
    give the common translation of every crossing label.
\end{proof}

\begin{lemma}[Counting with a restricted common translation]%
    \label{thm:peps_regular_twisted_boundary_count}
    \lean{TNLean.PEPS.regionBoundaryLabel_eq_iff,
        TNLean.PEPS.sum_boundaryRestriction_filtered_translations_eq_card,
        TNLean.PEPS.sum_regionBoundary_filtered_translations_eq_card}
    \leanok
    \uses{thm:peps_regular_twisted_synchronization}
    Let $I_R$ count the internal bonds and let $F\subseteq G$ be any
    allowed set of common translations. With crossing configurations
    $\mu,\nu$, the number of translated pairs of bond assignments is
    \begin{align}
        |G|^{I_R}\sum_{x\in F}\mathbf{1}_{\mu=x\nu}.
    \end{align}
\end{lemma}

\begin{proof}\leanok
    For fixed $x$, the second assignment determines the first: its crossing
    labels are fixed and every internal label is free, while the boundary conditions
    agree exactly when $\mu=x\nu$.
\end{proof}

\begin{lemma}[Weighted translations on the invariant boundary]%
    \label{thm:peps_regular_weighted_boundary_translation}
    \lean{TNLean.PEPS.toMatrix_regularBoundaryRepresentation_mul_projector,
        TNLean.PEPS.regularBoundaryWeightedTranslation,
        TNLean.PEPS.regularBoundaryWeightedTranslation_apply,
        TNLean.PEPS.regularBoundaryWeightedTranslation_mul_projector,
        TNLean.PEPS.regularBoundaryProjector_mul_weightedTranslation_mul_projector}
    \leanok
    \uses{def:peps_regular_boundary_state}
    For weights $a_x\in\mathbb C$, let
    $W_a=\sum_x a_x L_x^{\otimes b}$. Its matrix entries are
    $(W_a)_{\mu,\nu}=\sum_x a_x\mathbf{1}_{\mu=x\nu}$.
    If $P$ is the invariant-boundary projector, then
    \begin{align}
        W_aP=\left(\sum_xa_x\right)P,
        \qquad PW_aP=\left(\sum_xa_x\right)P.
    \end{align}
\end{lemma}

\begin{proof}\leanok
    Every common translation fixes every vector in the range of $P$;
    summing this identity with the specified weights and using $P^2=P$ yields both equalities.
\end{proof}

\begin{lemma}[Twisted-region cross Gram matrix]%
    \label{thm:peps_regular_twisted_region_cross_gram}
    \lean{TNLean.PEPS.sum_regularBoundary_filtered_translation_eq_weightedTranslation,
        TNLean.PEPS.regularTwistedOpenRegionMatrix_crossGram_of_connected,
        TNLean.PEPS.regularTwistedOpenRegionMatrix_projected_crossGram_of_connected}
    \leanok
    \uses{thm:peps_regular_twisted_synchronization, thm:peps_regular_weighted_boundary_translation}
    Under the hypotheses of
    Lemma~\ref{thm:peps_regular_twisted_synchronization}, put
    \begin{align}
        F_{u,w}=\{x\in G:u_ex=xw_e\text{ for every internal bond }e\},
        \qquad
        \kappa_R=\left(\prod_{v\in R}c_v\right)
        \frac{|G|^{I_R+1}}{|G|^{|R|}}.
    \end{align}
    The maps satisfy
    \begin{align}
        (T_R^u)^\dagger T_R^w
        =\frac{\kappa_R}{|G|}\sum_{x\in F_{u,w}}L_x^{\otimes b},
        \qquad
        P(T_R^u)^\dagger T_R^wP
        =\kappa_R\frac{|F_{u,w}|}{|G|}P.
    \end{align}
    These are auxiliary finite-region identities with identity operators
    on crossing bonds. The general torus cut of
    \cite[Theorem~6.9]{Schuch2010PEPS} remains separate
    \cite{gap:rmp_peps_quantum_double_g_isometry}.
\end{lemma}

\begin{proof}\leanok
    \uses{thm:peps_regular_twisted_boundary_count}
    By the expansion and synchronization lemmas, counting the free
    internal labels contributes $|G|^{I_R}$ and leaves the prefactor
    $\kappa_R/|G|$. On the invariant boundary each remaining translation
    acts identically, giving the projected formula.
\end{proof}

\begin{lemma}[Superpositions with a common boundary density]%
    \label{thm:peps_regular_closure_superposition_density}
    \lean{TNLean.PEPS.regularClosureSuperpositionMap_conjTranspose_mul,
        TNLean.PEPS.regularClosureSuperpositionMap,
        TNLean.PEPS.regularClosureSuperpositionFactor,
        TNLean.PEPS.regularClosureSuperpositionMap_mul_projector,
        TNLean.PEPS.regularClosureSuperpositionState,
        TNLean.PEPS.regularClosureSuperpositionState_eq_sum,
        TNLean.PEPS.regularClosureSuperpositionState_eq_sum_of_gram,
        TNLean.PEPS.exists_normalization_regularClosureSuperpositionState}
    \leanok
    \uses{thm:peps_regular_weighted_boundary_translation,
        thm:peps_regular_boundary_physical_realization}
    Let $b=n+1$, let $P$ be the regular boundary projector, and suppose
    specified cut maps satisfy
    \begin{align}
        A^\dagger A=c_AP,\quad c_A>0,\qquad
        PN_s^\dagger N_tP=k_{s,t}P.
    \end{align}
    For coefficients $\lambda_s$, put
    $B_\lambda=(\sum_s\lambda_sN_s)P$ and
    $K_\lambda=AB_\lambda^{\mathsf T}$. If $K_\lambda\ne0$, a positive
    scalar normalizes its coefficient vector, whose reduced density is
    \begin{align}
        \rho_A=
        \frac{1}{c_A|G|^n}APA^\dagger.
    \end{align}
    This density is independent of $\lambda$. It has rank $|G|^n$,
    satisfies $\rho_A^2=|G|^{-n}\rho_A$, and has entropy $n\log|G|$.
    This is a conditional matrix statement; identification of its Gram
    hypotheses with a PEPS contraction is separate.
\end{lemma}

\begin{proof}\leanok
    Expanding the mixed Gram matrices gives
    \begin{align}
        B_\lambda^\dagger B_\lambda=\kappa_\lambda P,\qquad
        \kappa_\lambda=\sum_{s,t}\overline{\lambda_s}\lambda_t k_{s,t}.
    \end{align}
    Nonvanishing of $K_\lambda$ implies $B_\lambda\ne0$, so positivity of
    its Gram matrix and its nonzero trace imply $\kappa_\lambda>0$.
    Normalizing $A$ and $B_\lambda$ by their positive square-root factors
    and applying the physical boundary density calculation completes the proof.
\end{proof}

\begin{lemma}[Common density for PEPS with complementary twists]%
    \label{thm:peps_regular_twisted_cut_entropy}
    \lean{TNLean.PEPS.regularTwistedSite_eq_of_trivial_on_incident,
        TNLean.PEPS.regularTwistedOpenRegionMatrix_eq_of_trivial_on_incident,
        TNLean.PEPS.regularClosureSuperpositionState_eq_twistedCut_sum,
        TNLean.PEPS.regularTwistedPhysicalCutSuperposition,
        TNLean.PEPS.exists_normalization_regularTwistedPhysicalCutSuperposition}
    \leanok
    \uses{thm:peps_regular_closure_superposition_density}
    Let every site tensor of a finite graph be regular $G$-isometric.
    Let $R$ be nonempty and connected, with $b=n+1$ crossing bonds.
    Consider a finite family of bond assignments $u^{(s)}$, all equal
    to the identity on edges incident to $R$. Suppose their common
    identity bonds contain a connected spanning subgraph of $R^c$.
    Let $\psi_s$ be the physical PEPS across the cut with
    assignment $u^{(s)}$.

    There is one $c_R>0$ such that the untwisted region map
    satisfies $T_R^\dagger T_R=c_RP$. Every nonzero superposition
    $\sum_s\lambda_s\psi_s$ has a positive normalization whose reduced
    density on $R$ is the same operator
    \begin{align}
        \rho_R=\frac{1}{c_R|G|^n}T_RPT_R^\dagger.
    \end{align}
    In particular its rank is $|G|^n$, its nonzero spectrum is flat,
    and its entropy is $n\log|G|$.
    This is an auxiliary finite-cut statement with twists confined to
    the complement, rather than the arbitrary torus cut of
    \cite[Theorem~6.9]{Schuch2010PEPS}
    \cite{gap:rmp_peps_quantum_double_g_isometry}.
\end{lemma}

\begin{proof}\leanok
    \uses{thm:peps_regular_twisted_region_cross_gram, thm:peps_open_region_cut}
    Identity operators on edges incident to $R$ leave its open
    map unchanged, and the connected-region calculation derives its positive
    Gram factor from the local isometries. The twisted-region calculation
    derives each projected complementary mixed Gram matrix as a scalar
    multiple of $P$. Exact contraction across the cut identifies the
    superposition with the matrix state of
    Lemma~\ref{thm:peps_regular_closure_superposition_density}, which then yields the common reduced density, rank, and entropy without assuming a Gram identity for a whole region.
\end{proof}

\begin{lemma}[Zero-order entropy of regular $G$-injective torus closures]%
    \label{thm:peps_regular_ginjective_simply_connected_entropy}
    \lean{TNLean.PEPS.IsGInjective.torusPhysicalCut_rank_zeroEntropy_of_isSimplyConnected}
    \leanok
    \uses{def:peps_renyi_entropy}
    Let both torus circumferences be at least three, and let $A$ be a
    $G$-injective site tensor with the regular virtual representation.
    Let $R$ be a contiguous block: its induced nearest-neighbour graph
    is connected, and its closed-cell realization is simply
    connected. Let $b$ be the number of bonds crossing its boundary.
    For any nonzero coherent superposition of closures with
    commuting pairs, let $M$ be its physical coefficient matrix across
    the cut and put
    \begin{align}
        \rho_R & =\frac{MM^\dagger}{\operatorname{tr}(MM^\dagger)}.
    \end{align}
    The denominator is positive, $\rho_R$ is a density matrix, and
    \begin{align}
        \operatorname{rank}\rho_R & =|G|^{b-1},                    &
        S_0(\rho_R)               & =\log\operatorname{rank}\rho_R
        =(b-1)\log|G|.
    \end{align}
    Trees, roots, boundary numbering, and complementary paths are not
    assumed. This is the regular native-torus conclusion of
    \cite[Corollary~6.10]{Schuch2010PEPS}, under the stated contiguous
    closed-cell geometry. It makes no flat-spectrum assertion and does
    not identify the closure span with a parent-Hamiltonian ground space
    \cite{gap:rmp_peps_quantum_double_g_isometry}.
\end{lemma}

\begin{proof}\leanok
    \uses{thm:peps_regular_g_injective_closure_rank_transport,
        thm:peps_torus_simply_connected_physical_density}
    The products of local inverse maps carry the original state to
    the canonical closure superposition. The products of the original
    site maps recover it. Thus the two coefficient matrices have equal rank. Numbering the canonical physical basis preserves its
    isometry and its cut rank, so the simply connected isometric
    density theorem gives the canonical rank $|G|^{b-1}$.
    A matrix and its row Gram matrix have equal rank; division by the
    positive trace preserves that rank. The definition of zero-order
    Rényi entropy gives the stated logarithm.
\end{proof}


\begin{theorem}[Simply connected rectangular cell regions]%
    \label{thm:peps_torus_rectangle_realization}
    \lean{TNLean.PEPS.isSimplyConnected_torusRegionRealization_rectangle}
    \leanok
    \uses{thm:peps_torus_region_realization}
    Let the torus periods be positive integers $w,h$. Let $a,b\geq0$ and
    $\ell,m>0$ be integers with $\ell<w$, $m<h$, $a+\ell\leq w$, and
    $b+m\leq h$, and let $R$ be the native rectangle with coordinate ranges
    $a\leq x<a+\ell$ and $b\leq y<b+m$.
    Then the closed-cell realization of $R$ is simply connected.
    This is auxiliary geometry for the stripe reduction in
    \cite[Theorem~6.7]{Schuch2010PEPS}, local source lines 1995--2011.
\end{theorem}

\begin{proof}\leanok
    The torus projection maps the planar rectangle
    $[a-1/2,a+\ell-1/2]\times[b-1/2,b+m-1/2]$ onto the closed-cell
    realization of $R$. To verify surjectivity, choose in each coordinate
    an integer center in the indicated range at distance at most $1/2$;
    conversely, every cell centered in $R$ lies in this planar rectangle.
    Since its side lengths are smaller than the periods, projection is
    injective there. The continuous bijection from this compact rectangle
    to its Hausdorff image is a homeomorphism. The planar rectangle is
    nonempty and convex, hence contractible and simply connected.
\end{proof}


\begin{theorem}[Unitary equivalence of closure states on two width-one stripes]%
    \label{thm:peps_torus_stripe_local_equivalence}
    \lean{TNLean.PEPS.mem_compl_torusStripeInterior_iff_coordinate_zero,
        TNLean.PEPS.IsGIsometric.torusClosureSuperpositionCut_stripe_local_equivalence}
    \leanok
    \uses{def:peps_normalized_pure_cut_matrix}
    Consider an original regular $G$-isometric tensor on a native
    torus whose two periods are at least three. Let $R$ be the
    nonwrapping rectangle of sites with both coordinates nonzero.
    Its complement is the union of two width-one stripes wrapping
    around the torus:
    \begin{align}
        S & =\{(x,y):x=0\lor y=0\}.
        \label{eq:peps_width_one_stripes}
    \end{align}
    For any two nonzero coherent superpositions of commuting torus
    closures, let $A$ and $B$ be their physical coefficient matrices
    across $R\mid S$. There exists a unitary $U$ on the physical
    factor of $S$ satisfying
    \begin{align}
        \widehat B & =\widehat A U^{\mathsf T}.
        \label{eq:peps_stripe_physical_unitary}
    \end{align}
    In this closure-vector specialization of
    \cite[Theorem~6.7]{Schuch2010PEPS} to regular native tori with periods
    at least three, $U$ transforms the first normalized physical state into
    the second while acting only on the stripes
    (\ref{eq:peps_width_one_stripes}); identifying the closure span with a
    parent-Hamiltonian ground space is handled separately.
\end{theorem}

\begin{proof}\leanok
    \uses{thm:peps_torus_rectangle_realization,
        thm:peps_torus_rectangle_cut_connectivity,
        thm:peps_torus_simply_connected_physical_local_equivalence}
    The coordinate-value bounds identify $R$ with the rectangle
    starting at $(1,1)$ with side lengths one less than the two periods.
    Its induced graph is connected, and
    \ref{thm:peps_torus_rectangle_realization} gives simple
    connectedness of its closed-cell realization, so applying
    \ref{thm:peps_torus_simply_connected_physical_local_equivalence}
    yields (\ref{eq:peps_stripe_physical_unitary}).
\end{proof}


\begin{lemma}[Closure seams avoid a strictly interior rectangle]%
    \label{thm:peps_torus_interior_rectangle_seam_exclusion}
    \lean{TNLean.PEPS.torusNonseamGraph_not_horizontal_wrap,
        TNLean.PEPS.torusNonseamGraph_not_vertical_wrap,
        TNLean.PEPS.torusRightEdge_not_incident_rectangle_of_wrap,
        TNLean.PEPS.torusUpEdge_not_incident_rectangle_of_wrap}
    \leanok
    \uses{def:peps_torus_rectangular_grid}
    Suppose both circumferences are at least three. Wrap bonds are absent
    from the ordinary rectangular grid. If a rectangle lies strictly
    between the first and last coordinate values in both directions,
    no wrap bond is incident to it.
\end{lemma}

\begin{proof}\leanok
    \uses{thm:peps_torus_rectangle_cut_connectivity}
    A wrap bond has coordinate values $0$ and $w-1$, or $0$ and $\ell-1$.
    Their difference cannot be one for a circumference at least three,
    and neither endpoint can lie in the stated interior rectangle.
\end{proof}

\begin{lemma}[Regular isometry in incident-edge coordinates]%
    \label{thm:peps_regular_torus_incident_isometry}
    \lean{TNLean.PEPS.torusIncidentCoordinates_injective,
        TNLean.PEPS.torusIncidentCoordinates,
        TNLean.PEPS.IsGIsometric.exists_torusIncidentSiteGram,
        TNLean.PEPS.IsGIsometric.isGIsometric_torusIncidentSite,
        TNLean.PEPS.torusClosureEdgeAssignment_eq_one_of_nonseam,
        TNLean.PEPS.torusClosureEdgeAssignment_eq_one_on_rectangle_incident}
    \leanok
    \uses{thm:peps_regular_torus_closure_cut_identification}
    Suppose both torus dimensions are at least three. A tensor which is
    $G$-isometric for the four-leg regular action remains $G$-isometric
    when its virtual coordinates are indexed by incident graph edges.
    Every closure assignment is the identity on the ordinary grid and
    on every edge incident to a strictly interior rectangle.
\end{lemma}

\begin{proof}\leanok
    \uses{thm:peps_torus_interior_rectangle_seam_exclusion}
    The four incident-edge labels determine the four virtual coordinates
    bijectively and intertwine simultaneous left multiplication. Thus
    their group-average Gram matrices agree. The claims about assignments
    follow from the exclusion of wrap bonds.
\end{proof}

\begin{lemma}[Entropy of closure superpositions on interior rectangles]%
    \label{thm:peps_regular_torus_interior_rectangle_entropy}
    \lean{TNLean.PEPS.torusClosureSuperpositionCut_eq_twistedCutSuperposition,
        TNLean.PEPS.torusClosureSuperpositionCut,
        TNLean.PEPS.exists_normalization_torusClosureSuperpositionCut_rectangle}
    \leanok
    Let $A$ be regular $G$-isometric on a torus of width and height at
    least three. Let $R$ be a coordinate rectangle of positive side
    lengths satisfying
    \begin{align}
        0<x_0,\quad x_0+x_1<w,\qquad
        0<y_0,\quad y_0+y_1<\ell.
    \end{align}
    Write $b$ for its number of crossing bonds and $P$ for the invariant
    boundary projector. There is one $c_R>0$ with
    $T_R^\dagger T_R=c_RP$. For any finite family of closure pairs
    $(g_s,h_s)$ and coefficients $\lambda_s$ such that
    $\sum_s\lambda_s\psi_{g_s,h_s}\ne0$, a positive normalization of this
    torus vector has reduced density on $R$
    \begin{align}
        \rho_R=\frac{1}{c_R|G|^{b-1}}T_RPT_R^\dagger.
    \end{align}
    Hence the normalized reduced density is independent of the
    superposition. It has rank $|G|^{b-1}$ and flat nonzero spectrum,
    and its von Neumann entropy is $(b-1)\log|G|$.
    This proves the strictly interior rectangular case of the entropy
    calculation in \cite[Theorem~6.9]{Schuch2010PEPS}. The present case
    also permits noncommuting closure pairs. The source's full statement
    concerns commuting sectors and arbitrary topologically trivial
    regions \cite{gap:rmp_peps_quantum_double_g_isometry}.
\end{lemma}

\begin{proof}\leanok
    \uses{thm:peps_torus_rectangle_cut_connectivity,
        thm:peps_regular_twisted_cut_entropy,
        thm:peps_regular_torus_incident_isometry}
    The coordinate bounds give connected $R$, a nonempty boundary, and
    connected $R^c$ after removing all wrap bonds, and the closure assignments
    are the identity on this complementary grid and on all edges incident
    to $R$. Rewriting the local regular isometry in incident-edge
    coordinates and identifying the native closure superposition with its
    twisted graph contraction allows
    Lemma~\ref{thm:peps_regular_twisted_cut_entropy} to apply: nonvanishing gives
    the positive normalization, and the same region map and constant apply
    to every choice of coefficients.
\end{proof}
